\chapter{Supplementary material}
\label{chap:appendix-a}

\section{Tail latency}

\begin{table}[htbp]
  \centering
  \caption{Build time percentiles for a \num{100}-page document.}
  \label{tab:percentiles}
  \small
  \begin{tabular}{@{}l S[table-format=2.2] S[table-format=2.2] S[table-format=2.2]@{}}
    \toprule
    \textbf{\sffamily Strategy} & {p50 (s)} & {p90 (s)} & {p95 (s)} \\
    \midrule
    Baseline    & 7.94 & 11.62 & 14.08 \\
    Cached      & 3.52 &  4.91 &  6.03 \\
    Incremental & 0.94 &  1.28 &  1.51 \\
    \bottomrule
  \end{tabular}
\end{table}

\section{Reproducing the measurements}

\begin{lstlisting}[language=bash,caption={Full reproduction sequence.},label={lst:repro}]
git clone https://github.com/<your-org>/<this-repo>.git
cd <this-repo>
make bench          # writes data/measurements.csv
latexmk -pdf main   # rebuilds this report from the new data
\end{lstlisting}

\begin{remark}
The figures in \cref{chap:results} read \texttt{data/measurements.csv}
at compile time, so no figure needs to be regenerated by hand after a
new benchmark run.
\end{remark}
